\section{授業内試験}
\label{b30_final_exam}

\subsection{授業内容}

\subsubsection{概要}
本コマでは，これまでに学んできた心理統計学の内容についての理解度を確認するための授業内試験を実施する。前期から後期にわたる授業で学んだ基本的な知識，重要な統計概念，データ分析の方法論，そして心理統計学の考え方の筋道について総合的に問う試験を通じて，一年間の学習の成果を測定する。試験の内容は，記述統計から推測統計，回帰分析，分散分析，そして近年重要性を増しているベイズ統計学までを含む幅広いものとなる。試験結果は成績評価の重要な要素となるため，これまでの学習内容を総復習し，心理統計学の基礎を確実に身につけておくことが望ましい。

\subsubsection{コマ主題細目}
\begin{description}
	\item[試験の概要と出題範囲] 授業内試験はマークシート形式で実施され，問題数は60〜70問程度である。出題範囲は前期から後期までの全授業内容を含み，特に以下の領域からの出題が予想される：(1)記述統計と確率の基礎概念，(2)推測統計と仮説検定の論理，(3)分散分析と実験計画法，(4)回帰分析と相関分析，(5)一般線形モデル，(6)ベイズ統計学の基礎，(7)統計ソフトウェア（RとJASP）の操作と結果の解釈。問題形式としては，基本的な用語や概念の理解を問う問題，計算問題，結果の解釈を問う問題，適切な分析手法の選択を問う問題などが含まれる。問題は基礎的な内容から応用的な内容まで幅広く出題され，単なる暗記ではなく，統計的思考力を測定することを目的としている。
	
	\item[学習の総括と重要概念の確認] この試験は，一年間の心理統計学の学習内容を総括する機会でもある。特に重要な概念として，(1)記述統計と推測統計の区別，(2)確率分布と標本分布の関係，(3)仮説検定の論理と$p$値の正しい解釈，(4)効果量と検定力の概念，(5)モデリングの考え方と一般線形モデルの枠組み，(6)頻度主義統計学とベイズ統計学の違いなどが挙げられる。これらの概念は心理学研究における統計分析の基礎となるものであり，卒業研究や大学院での研究活動においても重要となる。試験問題を通じてこれらの概念を再確認し，心理統計学の全体像を把握することが期待される。

	\item[実践的なデータ分析とソフトウェアの利用] 本授業では，理論的な内容の学習と並行して，RやJASPなどの統計ソフトウェアを用いた実践的なデータ分析のスキルも重視してきた。試験では，ソフトウェアの操作そのものではなく，分析結果の読み方や解釈，適切な分析手法の選択などについての理解を問う問題が含まれる。具体的には，出力結果からの情報抽出，結果の統計学的および実質的意味の解釈，分析結果に基づく結論の導出などが問われる。これらの能力は，心理学研究の実践において必須のスキルであり，試験を通じてこれらの能力が身についているかを確認する。

	
\end{description}
\subsubsection{キーワード}
\begin{itemize}
	\item 心理統計学の総括
	\item 統計的思考力
	\item 分析結果の解釈
	\item 頻度主義統計学とベイズ統計学
\end{itemize}
\subsection{授業情報}
\paragraph{コマの展開方法}
授業内試験
\paragraph{標準シラバスにおける位置づけ}
科目番号5；心理学統計法；総括
\subsubsection{予習・復習課題}
\paragraph{予習}
本コマは授業内試験であるため，以下の内容について十分な準備をして臨むこと：

\begin{enumerate}
\item 前期および後期の授業内容全体を体系的に復習し，重要な概念と用語を理解する
\item 配布資料や教科書の内容を再確認し，特に重要な統計手法とその適用条件について整理する
\item これまでに行った実習や課題の内容を振り返り，分析結果の解釈方法を確認する
\item 試験の詳細（持ち込み可能な資料の範囲など）については，別途配布されるお知らせを参照すること
\item 過去の小テストや課題で理解が不十分だった部分を特に重点的に復習する
\end{enumerate}

\paragraph{復習}
試験終了後は，問題内容を振り返り，不明点があれば次回の授業（総括回）で質問できるよう準備しておくとよい。心理統計学の基礎を確実に身につけることは，今後の心理学研究において重要な基盤となるため，試験を通じて得られた知識の定着を図ることが重要である。また，試験の結果を踏まえて，自分の弱点や苦手分野を特定し，今後の学習計画に活かすことも有益である。